\section{Uniform averaging of arithmetic signals}\label{sec:applications}

The distinction between a single rule that works at every translate and a rule
chosen separately at each translate is essential.  The kernels constructed
above are independent of the center of the averaging interval.  Their
application therefore gives uniform control of an entire almost periodic
signal, with the same positive weights at every center.

\begin{theorem}[A uniform positive averaging rule]\label{thm:app-uniform}
Let $\Gamma\subset\mathbb R\setminus\{0\}$ be a locally finite symmetric set,
and let
\[
 f(t)=b+\sum_{\gamma\in\Gamma}c_\gamma e^{i\gamma t},
 \qquad \mathcal A=\sum_{\gamma\in\Gamma}\|c_\gamma\|<\infty,
\]
where the coefficients take values in a finite-dimensional normed space.
Suppose the low-frequency annihilation construction, followed by its uniform
convolution factor, gives a probability measure $\mu_B$ with
\[
 \operatorname{supp}\mu_B\subset[-H_B,H_B],\qquad
 \widehat\mu_B(\gamma)=0\quad(0<|\gamma|\le B),\qquad
 \sup_{|\omega|\ge B}|\widehat\mu_B(\omega)|\le e^{-c_0B\log B}
\]
for $B\ge2$, where $\widehat\mu_B(\omega)=\int e^{i\omega u}\,d\mu_B(u)$
and $c_0=-\log(\sin 1)>0$ is permissible.  Then
\begin{equation}\label{eq:app-uniform-error}
 \sup_{t\in\mathbb R}
 \left\|\int f(t+u)\,d\mu_B(u)-b\right\|
 \le \mathcal A e^{-c_0B\log B}.
\end{equation}
In particular, if $H_B\le C\log B$ for $B\ge B_0$, then, as
$\varepsilon\downarrow0$, a single positive averaging rule of half-width
\[
 H=O\!\left(\log\log(\mathcal A/\varepsilon)\right)
\]
achieves error at most $\varepsilon$ at every center.  This last asymptotic
statement is understood for $\mathcal A>0$; if $\mathcal A=0$, no averaging
is necessary.
\end{theorem}

\begin{proof}
Absolute uniform convergence permits interchange of the Fourier series and
integration.  The difference from $b$ is
\[
 \sum_{|\gamma|>B}c_\gamma e^{i\gamma t}\widehat\mu_B(\gamma).
\]
Taking its norm proves \eqref{eq:app-uniform-error}, uniformly in $t$.
For the accuracy assertion put $L=\log(\mathcal A/\varepsilon)$ and take,
for example, $B=2L/(c_0\log L)$ once $L$ is sufficiently large.  Then
$c_0B\log B\ge L$ and
$\log B=\log L-\log\log L+O(1)$.
\end{proof}

The assertion is about uniform approximation by a common convolution kernel.
It does not assert that every procedure for balancing values in an interval
needs a growing interval.  The earlier ProveIt convex-hull theorem \cite{priorbalance} already
allows center-dependent positive atomic measures to balance the finite vector
of ramp errors exactly in a fixed-width logarithmic interval.  These measures
can change when the center changes.  There is no conflict with the present
translation-independent construction.  Likewise, the Jensen lower bound above
establishes optimal order of support for \emph{exact low-frequency
annihilation}; it is not, without an additional argument, a lower bound for
all approximate averaging procedures.

\subsection{The ramp sum and its sharper remainder}

Fix $q\ge1$, let $d=\varphi(q)$, and let $a$ range over the reduced residue
classes modulo $q$ (the group has one element when $q=1$).  Define
\begin{align}
 S_a(x)&=\sum_{\substack{n\le x\\ n\equiv a\pmod q}}
              \Lambda(n)\left(\frac1n-\frac1x\right),\notag\\
 D_a(s)&=\sum_{n\equiv a\pmod q}\frac{\Lambda(n)}{n^s}
       =-\frac1d\sum_{\chi\bmod q}\overline{\chi(a)}
                         \frac{L'}{L}(s,\chi),\qquad \Re s>1,\notag\\
 B_a&=\lim_{s\to1}\left(D_a(s)-\frac{1}{d(s-1)}\right),\notag\\
 R_a(x)&=S_a(x)-\frac{\log x}{d}-B_a+\frac1d,
       \qquad F_a(v)=e^{v/2}R_a(e^v).\label{eq:app-ramp-def}
\end{align}
The constant $B_a$ exists by the classical nonvanishing of Dirichlet
$L$-functions at $1$.  These are the same ramp normalizations as in the
previous simultaneous-balancing report \cite{priorbalance}.  For a character $\chi$ modulo $q$,
let $\chi^*$ denote its primitive inducing character, and let
$m_\chi(\gamma)$ denote the multiplicity of the zero of $L(s,\chi^*)$ at
$s=1/2+i\gamma$, or zero if there is no such zero.  Set
\begin{align}
 b_a&=\frac4d\sum_{\chi\bmod q}\overline{\chi(a)}m_\chi(0),\notag\\
 c_a(\gamma)&=
 \frac{\sum_{\chi\bmod q}\overline{\chi(a)}m_\chi(\gamma)}
 {d(1/4+\gamma^2)}\qquad(\gamma\ne0),\notag\\
 \mathcal A_q&=\frac1d\sum_{\chi\bmod q}
                  \sum_{\gamma\ne0}\frac{m_\chi(\gamma)}{1/4+\gamma^2}.
 \label{eq:app-ramp-coeff}
\end{align}
The inner sums in the last expression use distinct signed ordinates.  In
particular $\sum_\gamma|c_a(\gamma)|\le\mathcal A_q<\infty$ for every
$a$.  Character conjugation makes $b_a$ and $F_a$ real and gives
$c_a(-\gamma)=\overline{c_a(\gamma)}$.

\begin{lemma}[A remainder of order $v e^{-v/2}$]\label{lem:app-ramp}
Assume Dirichlet GRH.  For each fixed $q$, uniformly for $v\ge0$,
\begin{equation}\label{eq:app-ramp-expansion}
 F_a(v)=b_a+\sum_{\gamma\ne0}c_a(\gamma)e^{i\gamma v}
                 +O_q\bigl((1+v)e^{-v/2}\bigr).
\end{equation}
The Fourier series is absolutely and uniformly convergent.  For a primitive
character the error has the more precise form
\[
 (r_\chi v+r_\chi-d_\chi)e^{-v/2}+O_\chi(e^{-3v/4}),
\]
where $-L'/L(s,\chi)=-r_\chi/s+d_\chi+O(s)$ at $s=0$.
\end{lemma}

\begin{proof}
First work with a primitive character $\chi$, allowing the conductor-one
character.  Write $D_\chi=-L'/L$, let $\delta_\chi=1$ for the
conductor-one character and $0$ otherwise, and use the analogous scalar
definitions $S_\chi,B_\chi,R_\chi$ with main term
$\delta_\chi\log x$.  The standard explicit-formula contour method \cite{Davenport} begins with the absolutely convergent
vertical integral
\[
 S_\chi(x)=\frac1{2\pi i}\int_{2-i\infty}^{2+i\infty}
       \frac{D_\chi(s)x^{s-1}}{s(s-1)}\,ds.
\]
Shift the contour to $\Re s=-1/4$.  There are no trivial zeros on this
line.  The functional equation and the absolutely convergent Euler series
on $\Re(1-s)=5/4$ give
\[
 D_\chi(-1/4+it)=O_\chi(\log(2+|t|)).
\]
The gamma terms obey the same estimate by Stirling's formula; on a compact
part of the line boundedness follows because it meets no zero or pole.
Consequently the new vertical integral is absolutely convergent and
$O_\chi(x^{-5/4})$, uniformly for $x\ge1$.

For clarity, the contour shift can be performed along heights that avoid
the zeros of the fixed finite family.  The standard local zero estimate
$N_\chi(t+1)-N_\chi(t)=O_q(\log(2+|t|))$ allows heights whose distance
from every nearby ordinate is at least a sufficiently small constant times
$1/\log(2+|t|)$.  The local partial-fraction formula then bounds the
logarithmic derivatives by $O_q(\log^2(2+|t|))$ along the horizontal
segments.  Their integrals tend to zero because the kernel contains
$s(s-1)$ in the denominator.  This argument works separately at the top
and bottom of the contour.  The residues at the critical zeros converge
absolutely after the shift, since
$N_\chi(U)=O_\chi(U\log(U+2))$.

At $s=1$ the residue is
$\delta_\chi\log x+B_\chi-\delta_\chi$.  A nontrivial zero $\rho$
of multiplicity $m$ contributes
\[
 \frac{m x^{\rho-1}}{\rho(1-\rho)}.
\]
At $s=0$, expand
\[
 D_\chi(s)=-\frac{r_\chi}{s}+d_\chi+O(s),\qquad
 \frac1{s(s-1)}=-\frac1s-1+O(s).
\]
The residue is therefore
$x^{-1}(r_\chi\log x+r_\chi-d_\chi)$.
The integer $r_\chi$ is $1$ for an even primitive nonprincipal character
and $0$ for an odd primitive character or the conductor-one character.
These values follow from the functional equation and nonvanishing at $1$.
No other trivial zero lies between the two vertical lines.  Under GRH,
$\rho(1-\rho)=1/4+\gamma^2$; multiplying by $\sqrt x$ proves the stated
primitive expansion.

It remains to justify that removing Euler factors does not spoil this
remainder.  If $\chi$ is induced by $\chi^*$, each omitted prime $p\mid q$
with $c=\chi^*(p)\ne0$ subtracts from the ramp sum
\[
 H_p(x)=(\log p)\sum_{1\le k\le\log x/\log p}
                         c^k(p^{-k}-x^{-1}).
\]
Its limiting constant is
$H_p(\infty)=(\log p)c/(p-c)$, and
\[
 H_p(x)-H_p(\infty)
 =-(\log p)\sum_{k>\log x/\log p}(c/p)^k
   -\frac{\log p}{x}\sum_{1\le k\le\log x/\log p}c^k
 =O_q\bigl((1+\log x)/x\bigr).
\]
The constant $B_\chi$ changes by precisely $-H_p(\infty)$, so the
corresponding error in $R_\chi$ is the negative of this difference.
There are finitely many omitted primes.  Finally apply character
orthogonality to pass from the scalar character sums to the residue-class
sums in \eqref{eq:app-ramp-def}.  The zero count proves absolute
summability in \eqref{eq:app-ramp-coeff}, completing the proof.
\end{proof}

\begin{theorem}[Simultaneous averaging on a polylogarithmic range of scales]
\label{thm:app-arithmetic}
Assume Dirichlet GRH and fix $q$.  There is a constant $C_q>0$ such that,
for all sufficiently large $X$, one can choose a single positive probability
measure $\mu_{\log X}$ supported on
\[
 [-C_q\log\log X,\ C_q\log\log X]
\]
for which, simultaneously for all reduced classes $a$ modulo $q$,
\begin{equation}\label{eq:app-polylog-error}
 \int e^{u/2}R_a(Xe^u)\,d\mu_{\log X}(u)
 =\frac{b_a}{\sqrt X}
   +O_q\!\left(\frac{(\log X)^{1+C_q/2}}{X}\right).
\end{equation}
Every scale used lies between
$X/(\log X)^{C_q}$ and $X(\log X)^{C_q}$.
For each bandwidth $B$, the kernel $\mu_B$ is independent of the center
$\log X$; the displayed rule uses the permissible choice $B=\log X$.
For $q=1$ the central bias is zero.
\end{theorem}

\begin{proof}
Apply the low-frequency construction to the union of the signed noncentral
zero ordinates of the primitive inducing characters modulo $q$.  The finite
family has $O_q(\log(B+2))$ zeros in each unit interval, with multiplicity,
so the construction provides $H_B\le C_q\log B$ after increasing $C_q$
and the lower threshold for $B$.  Frequencies below $1$, if present, form a
fixed finite set away from zero and can be included at a fixed additional
support cost.  The additional uniform convolutions give
$e^{-c_0B\log B}$ attenuation above $B$.

Put $T=\log X$, $B=T$ and $H=H_T$.  By
\eqref{eq:app-ramp-expansion} and Theorem~\ref{thm:app-uniform},
\[
 \left|\int F_a(T+u)\,d\mu_T(u)-b_a\right|
 \le \mathcal A_q e^{-c_0T\log T}
       +O_q\bigl((1+T+H)e^{-(T-H)/2}\bigr).
\]
The left side of \eqref{eq:app-polylog-error} is $X^{-1/2}$ times
$\int F_a(T+u)\,d\mu_T(u)$.  Since $H\le C_q\log T$, its arithmetic
remainder is $O_q(X^{-1}T^{1+C_q/2})$.  The spectral remainder is smaller
than $X^{-1}$ for all sufficiently large $T$.  The claimed scale interval
follows by exponentiating the support bound.  The value $\zeta(1/2)$ is
nonzero, so the conductor-one case has no central term.
\end{proof}

The factor $e^{u/2}$ in \eqref{eq:app-polylog-error} is part of the
normalization and cannot be omitted: the Fourier expansion belongs to
$F_a(v)=e^{v/2}R_a(e^v)$, not to $R_a(e^v)$ itself.  The conclusion
improves the support order from the earlier $O_q((\log\log X)^2)$
construction to $O_q(\log\log X)$, and replaces its
$X^{-1+\varepsilon}$ remainder by $X^{-1}$ times a fixed power of
$\log X$.  The power is not asserted to be optimal.  The bias term is
retained because GRH does not assert central nonvanishing of every
Dirichlet $L$-function.

\subsection{Positive atomic quadrature}

\begin{corollary}[At most $d+1$ actual scales]\label{cor:app-atoms}
Let $f:[-H,H]\to\mathbb R^d$ be continuous and let $\mu$ be a probability
measure supported on this interval.  There are $k\le d+1$ points
$u_1,\ldots,u_k\in[-H,H]$ and weights $w_j\ge0$ with
$\sum_jw_j=1$ such that
\[
 \sum_{j=1}^k w_j f(u_j)=\int f(u)\,d\mu(u).
\]
Thus the average in Theorem~\ref{thm:app-arithmetic} can, at each $X$, be
realized with at most $\varphi(q)+1$ prime scales, with exactly the same
error bound.
\end{corollary}

\begin{proof}
The image of the compact interval is compact, and its convex hull in
$\mathbb R^d$ is compact.  The integral belongs to this hull, as follows
by approximating the integral by positive finite sums, or by the separating
hyperplane theorem.  Carath\'eodory's theorem yields the representation.
Apply it to the continuous vector
$(e^{u/2}R_a(Xe^u))_a$.  Each ramp summand enters continuously, with value
zero when $n=x$, so continuity also holds at integer scales.
\end{proof}

The atomic locations and weights in this corollary can depend on $X$.
The underlying continuous convolution kernel is the object that is
independent of the center.  A fixed atomic quadrature is not asserted to
reproduce the entire infinite Fourier series at every translate.

\section{Robustness and limits of numerical certification}
\label{sec:app-robustness}

\begin{lemma}[Leakage from an approximate frequency]\label{lem:app-leakage}
Let $\mu$ be a probability measure supported on $[-A,A]$, and write
$\phi(\omega)=\int e^{i\omega t}\,d\mu(t)$.  Then
\[
 |\phi(\omega)-\phi(\lambda)|\le A|\omega-\lambda|.
\]
Suppose a convolution kernel is a product of probability-measure factors,
and one factor of half-width $A$ vanishes at an approximation $\lambda$ to
a true frequency $\gamma$.  Its full Fourier transform at $\gamma$ has
modulus at most $A|\gamma-\lambda|$.  If $m$ separate factors vanish at
$\lambda_1,\ldots,\lambda_m$ and have half-widths $A_1,\ldots,A_m$, the
stronger bound is
\[
 |\widehat\mu(\gamma)|\le
       \prod_{r=1}^m\min\{1,A_r|\gamma-\lambda_r|\}.
\]
\end{lemma}

\begin{proof}
Use $|e^{i\omega t}-e^{i\lambda t}|\le |\omega-\lambda|\,|t|$ and
integrate.  The remaining factors of a characteristic function have
modulus at most one.  Multiplying the individual estimates proves the
last assertion.
\end{proof}

Consequently, when every low frequency has an approximation of error at
most $\eta$ in a factor of half-width at most $A$, the uniform signal
bound becomes
\[
 \sup_t\left\|\int f(t+u)\,d\mu(u)-b\right\|
 \le \mathcal A\bigl(A\eta+e^{-c_0B\log B}\bigr),
\]
provided all true frequencies up to $B$ have been covered.  A known
multiplicity distributed among distinct factors gives the higher-order
leakage bound in the lemma.  The constant here is the width of a single
factor, rather than the width of the whole convolution.  For approximate
frequencies known only through intervals, one must also certify coverage
near the cutoff $B$; a slightly enlarged input band handles that boundary.

The positive-density construction is stable under perturbed nodes.  If
nodes initially separated by $\delta$ move by at most $\eta<\delta/2$,
their new separation is at least $\delta-2\eta$; rebuilding the factor with
$A(\delta-2\eta)\ge4\pi$ preserves the same matrix and positivity
estimates.  The positive and negative approximate nodes should be made
exactly symmetric.  One must also preserve the separation from zero.

There is a direct residual certificate for the finite linear solves.  Let
$M a=-b$ be the interpolation system for a symmetric separated set, so
$\|M-I\|_1\le1/12$, and let $\widetilde a$ be a symmetric real numerical
candidate.  If $r=M\widetilde a+b$, then
\[
 \|\widetilde a-a\|_1\le\frac{12}{11}\|r\|_1.
\]
If $\|\widetilde a\|_1<1$, the corresponding unnormalized density
$\,g_A(t)(1+\sum_j\widetilde a_j e^{-i\lambda_jt})$ is positive, its
normalizing integral is at least $1-\|\widetilde a\|_1$, and the
normalized Fourier residual at each prescribed node is at most
\[
 \frac{\|r\|_\infty}{1-\|\widetilde a\|_1}.
\]
These inequalities give a route from interval-certified nodes and matrix
residuals to a rigorous finite numerical kernel.  Ordinary floating-point
checks alone are illustrations of the construction: they do not certify a
complete list of $L$-function zeros, their multiplicities, or the analytic
theorems.

\subsection{Independent checks of the analytic argument}

The following points are useful when assessing the theorem's scope.
Repeated zeros are counted when coloring a multiset if the desired
annihilation order is their multiplicity.  Separate copies of the same
frequency must then be placed in different colors, and convolution adds
their vanishing orders.  If only ordinary annihilation is required, each
distinct frequency needs only one copy.  For a family of $L$-functions,
using the union with summed multiplicities is a sufficient, possibly
stronger requirement; no disjointness of their zero sets is used.

For a zero-counting lower bound based only on distinct frequencies, write
$N=\sum m$, $M_2=\sum m^2$, and $D=\sum1$ over distinct zeros in the
range.  The elementary inequality
$m-1\le(m^2-m)/2$ gives
\[
 D\ge\frac{3N-M_2}{2}.
\]
Thus any separately established estimate
$M_2\le(4/3+o(1))N$ implies $D\ge(5/6-o(1))N$.  Applying Jensen with
this count yields the corresponding $5/12$ leading lower coefficient for
the zeta zero set.  This elementary implication is independent of the
input theorem supplying the bound on $M_2$.

Finally, the tail estimate for the added uniforms uses the whole half-line:
$\sin x/x$ decreases on $[1,\pi]$, since the derivative's numerator
$x\cos x-\sin x$ has derivative $-x\sin x<0$ on $(0,\pi)$ and vanishes
at zero.  For $x\ge\pi$, $|\sin x/x|\le1/\pi<\sin1$.  Hence
$\sup_{|x|\ge1}|\operatorname{sinc}x|\le\sin1$ as required; restricting
the check to the first lobe would not suffice.
